% Draft supplementary-material insert registered at the v1.8.2 gate.
% The accepted source-v1 paper remains frozen.

\subsection{Scheduled meta-orders across operational and calendar time}
\label{app:scheduled-meta-orders}

A numerical meta-order is represented as a schedule of child market-order
events, not as one enlarged instantaneous density perturbation.  For child
$r$ with signed quantity $\epsilon q_r$ and declared operational index $n_r$,
the event map first consumes the opposing density from the current reaction
front outwards.  The ordinary uniform operational update is then applied from
$n_r-1$ to $n_r$.  Thus every child acts on the state left by the preceding
children and intervening diffusion, cancellation, source and coupling steps.

The control path and each scheduled-event path use the same initial state and
the same supplied operational innovations.  Both books are completed on the
uniform operational grid before either book-specific clock is introduced.
For calendar time $t$, book $j$ is observed only through its declared
previous-refresh subordination,
\begin{equation}
  P_j(t)=p_j\!\left(U_j(t)\right),
  \qquad
  U_j(t)=\max\{u_{j,m}:t_{j,m}\leq t\}.
\end{equation}
No interpolation and no state update on a nonuniform calendar grid are used.
Operational execution effects and observation-clock effects can therefore be
reported separately for exactly the same completed path.

The production experiment compares schedules with equal total signed volume
but different first-to-last-child horizons.  Total volume divided by this
horizon is a schedule-intensity proxy, not a participation rate: a true
participation rate would also require a defined background stream of market
order volume.  Post-completion own-impact relaxation, cross-impact catch-up and
long-lag own/cross convergence are reported as different observables because
coupling can transmit impact between children during a slow execution.

This construction inherits the translation-mode numerical source in
Equation~\eqref{eq:numerical-translation-source}.  No independent regularised
boundary width is added for a meta-order.  The resolved front changes after
each child, and subsequent coupling uses that current front profile.  The
bounded source in the accepted analytical construction remains a weak-moment
closure; the translation-mode numerical source conforms to its reduced
front-displacement law after projection and is not claimed to be pointwise
identical.
